Pour chacune des affirmations suivantes, indiquer si elle est vraie ou fausse, en justifiant votre réponse.

	Dans tout l'exercice, on se place dans un repère orthonormé.

	\begin{enumerate}[itemsep=5mm]
		\item Soit $f$ la fonction définie sur $\R$ par \[f(x)=2(x+1)(x-3)(x-5).\]

		\textbf{Affirmation 1 :} La fonction $f$ est négative sur l'intervalle $[-1 ; 3]$.

		\item Soit $h$ la fonction définie sur $\R$ par
		\[h(x)=-3 x^{2}+12 x-4.\]

		On note $\mathcal{C}_{h}$ sa courbe représentative.

		\begin{enumerate}
			\item \textbf{Affirmation 2 :} L'équation réduite de la tangente à $\mathcal{C}_{h}$ au point d'abscisse 1 est :
			\[y=6 x-1\]

			\item \textbf{Affirmation 3 :} La fonction $h$ est croissante sur l'intervalle $[2 ; +\infty[$.
		\end{enumerate}

		\item On considère la fonction polynôme de degré 2 définie sur $\R$ par :
		\[g(x)=\left(x-x_{1}\right)\left(x-x_{2}\right)\]

		où $x_{1}$ et $x_{2}$ sont des réels.

		On note $\mathcal{C}_{g}$ sa courbe représentative.

		On sait que :

		\begin{itemize}[label={\textbullet~}]
			\item 2 est une racine de $g$;
			\item la droite d'équation $x=2,5$ est l'axe de symétrie de $\mathcal{C}_{g}$.
		\end{itemize}

		\textbf{Affirmation 4 :} $g(0)=6$.

	\end{enumerate}
